% Nạp tệp này sau Chương 7 trong file chính DoAn.tex.
\documentclass[../DoAn.tex]{subfiles}

\begin{document}
\appendix
\chapter{Quy trình xây dựng và gán nhãn QCDT}
\label{appendix:qcdt-construction-labeling}

QCDT là tập hỏi đáp nội bộ được tác giả đồ án xây dựng và gán nhãn thủ công từ
Quyết định số 5445/QĐ-ĐHBK ban hành Quy chế đào tạo của Đại học Bách khoa Hà
Nội~\cite{hust2025trainingregulation}. Tập gồm 40 câu hỏi trên một tài liệu,
được chọn có chủ đích để bao phủ các điều khoản từ Điều 1 đến Điều 48 và các
nhóm nội dung chính: khái niệm, thời lượng và tín chỉ, đánh giá học phần, đăng
ký học tập, điều kiện tốt nghiệp, đào tạo sau đại học và quy định hiệu lực. Việc
chọn mẫu không ngẫu nhiên; mục tiêu là tạo các tình huống tra cứu có căn cứ trong
miền quy chế tiếng Việt.

Mỗi câu hỏi có một đáp án chuẩn \texttt{gold\_answer}, một hoặc nhiều trang
chuẩn \texttt{gold\_pages}, một hoặc nhiều mục/điều chuẩn \texttt{gold\_sections}
và nhãn loại bằng chứng \texttt{evidence\_type}. Trong 40 câu, có 20 câu
\texttt{factoid}, 2 câu \texttt{definition}, 10 câu \texttt{policy}, 4 câu
\texttt{comparison}, 2 câu \texttt{procedural} và 2 câu
\texttt{ambiguous\_or\_insufficient}. Về loại bằng chứng, tập có 25 câu dựa
trên đoạn văn, 10 câu dựa trên danh sách điều kiện, 4 câu dựa trên bảng và 1 câu
so sánh. Có 36 câu gắn với một trang chuẩn, 4 câu có hai trang chuẩn; 39 câu gắn
với một mục/điều và câu q39 gắn với hai điều để kiểm tra tổng hợp thông tin.

Đáp án chuẩn được viết bằng cách đọc trực tiếp phần quy chế đã gắn nhãn và giữ
lại các điều kiện, ngưỡng số hoặc quan hệ so sánh cần thiết để trả lời câu hỏi.
Mỗi câu được gán một đáp án chuẩn dạng chuỗi; tập hiện chưa có các phương án đáp
án tương đương do người gán nhãn thứ hai xây dựng. Việc gán nhãn độc lập bởi
người thứ hai và tính độ đồng thuận giữa người gán nhãn chưa được thực hiện, nên
QCDT cần được đọc như một tập đánh giá nội bộ có kiểm soát thay vì một chuẩn
đánh giá đã được xác nhận liên chủ thể.

Một số nhãn chuẩn tiêu biểu được trình bày tại
Bảng~\ref{tab:qcdt-label-examples}.

\begin{table}[H]
\centering
\caption{Ví dụ nhãn chuẩn trong tập QCDT}
\label{tab:qcdt-label-examples}

\scriptsize
\begin{tabularx}{\textwidth}{|
>{\centering\arraybackslash}p{0.09\textwidth}|
>{\raggedright\arraybackslash}X|
>{\raggedright\arraybackslash}p{0.18\textwidth}|
>{\raggedright\arraybackslash}p{0.22\textwidth}|
}
\hline
\textbf{Mã} & \textbf{Câu hỏi rút gọn} & \textbf{Loại bằng chứng} & \textbf{Trang và mục/điều chuẩn} \\
\hline
q01 & Phạm vi áp dụng và các văn bằng được cấp theo quy chế. & Đoạn văn & Trang 5, Điều 1 \\
\hline
q12 & Cộng hoặc trừ điểm quá trình theo số lần vắng mặt. & Bảng & Trang 9, Điều 5 \\
\hline
q20 & Điều kiện mở lớp học phần có giờ lên lớp và các ngoại lệ. & Danh sách điều kiện & Trang 12--13, Điều 10 \\
\hline
q23 & Điều kiện xét công nhận tốt nghiệp đại học. & Danh sách điều kiện & Trang 15, Điều 14 \\
\hline
q39 & So sánh số tín chỉ tối đa của sinh viên đại học và học viên thạc sĩ. & So sánh nhiều mục & Trang 12 và 23, Điều 10 và Điều 27 \\
\hline
\end{tabularx}
\end{table}

Trong bộ đánh giá hiện tại, \texttt{Answer Match} được tính đúng nếu câu trả
lời chứa đáp án chuẩn hoặc có Token F1 với \texttt{gold\_answer} không nhỏ hơn
0,45. \texttt{Evidence Match} được tính đúng nếu ít nhất một kết quả truy xuất
hoặc dẫn chứng nguồn khớp với một trang chuẩn hoặc chứa tên một mục/điều chuẩn.
QCDT chưa gán nhãn \texttt{expected\_chunk\_ids} hay văn bản bằng chứng chuẩn
cho từng câu; vì vậy, \texttt{Evidence Match} ở đây là độ đo khớp vị trí bằng
chứng ở mức trang/mục, không xác nhận khớp đúng từng khối nội dung và cũng không
đòi hỏi thu hồi đủ mọi bằng chứng trong các câu nhiều nguồn. Quy tắc này giải
thích vì sao \texttt{Evidence Match} có thể đạt 1,000 trong khi
\texttt{Answer Match} chỉ đạt 0,725.

\end{document}
